% GENERATED FILE. Edit canonical lessons, then run npm run generate:books.
% canonical-source-hash: fnv1a64:c1eb375418b05701
% canonical-lessons: KA-C473-parisara, KA-C473-malinya, KA-C473-kalusitagolisu, KA-C473-masi, KA-C473-ingala

\chapter{Our Environment}
\label{ch:ka-our-environment}
\hlchaptermodality{473}

\begin{chapteropening}
\textbf{By the end of this chapter:} \emph{I can say the environment, pollution, to pollute, soot and carbon.}

Complete the last lesson of chapter 473: I can say the environment, pollution, to pollute, soot and carbon.
\end{chapteropening}

\section[\texorpdfstring{\kn{ಪರಿಸರ} (parisara)}{ಪರಿಸರ (parisara)}]{\kn{ಪರಿಸರ} (parisara) — the environment; surroundings}

\label{lesson:KA-C473-parisara}



\begin{quote}
Before the new one: say the Kannada for shelter, then the Kannada for a helpline.
\end{quote}

\subsection*{You'll want to know: \kn{ಪರಿಸರ}}
\textbf{\kn{ಪರಿಸರ}} — \emph{parisara} — \textquotedblleft{}the environment; surroundings\textquotedblright{}.

\begin{quote}
\kn{ಪ್ಲಾಸ್ಟಿಕ್} \kn{ಪರಿಸರಕ್ಕೆ} \kn{ಒಳ್ಳೆಯದಲ್ಲ}.
\end{quote}

\begin{quote}
\emph{plāsṭik parisarakke oḷḷeyadalla.}
\end{quote}

\begin{quote}
Plastic is not good for the environment.
\end{quote}

\kn{ಪರಿಸರ} \kn{ದಿನ}, environment day.

\begin{cousinweb}
\textbf{\kn{ಪರಿಸರ}} (\emph{parisara}) is borrowed from Sanskrit \emph{parisara}, \textquotedblleft{}adjacent, lying around\textquotedblright{} (Monier-Williams); modern Kannada uses it for the environment.
\end{cousinweb}

\begin{grammarlens}[title={What you've built}]
One more word for getting around.
\end{grammarlens}

\subsection*{Your turn}
\begin{itemize}
\raggedright
  \item \emph{Say it:} \emph{parisara}
  \item \emph{Say it:} \emph{parisara}, once more
  \item \emph{Say it:} \emph{plāsṭik parisarakke oḷḷeyadalla.} — Plastic is not good for the environment.
  \item \emph{Recall:} say the Kannada for shelter, then the Kannada for a helpline, then say \emph{parisara} again
\end{itemize}

\begin{tcolorbox}[breakable,colback=teal!4,colframe=teal!35!black,title={Before you move on}]
What does \kn{ಪರಿಸರ} mean? (\textquotedblleft{}The environment; surroundings\textquotedblright{}.) Say it once more.
\end{tcolorbox}
\section[\texorpdfstring{\kn{ಮಾಲಿನ್ಯ} (mālinya)}{ಮಾಲಿನ್ಯ (mālinya)}]{\kn{ಮಾಲಿನ್ಯ} (mālinya) — pollution}

\label{lesson:KA-C473-malinya}



\begin{quote}
Before the new one: say the Kannada for a helpline, then the Kannada for the environment.
\end{quote}

\subsection*{You'll want to know: \kn{ಮಾಲಿನ್ಯ}}
\textbf{\kn{ಮಾಲಿನ್ಯ}} — \emph{mālinya} — \textquotedblleft{}pollution\textquotedblright{}.

\begin{quote}
\kn{ನಗರದಲ್ಲಿ} \kn{ವಾಹನಗಳಿಂದ} \kn{ಗಾಳಿಯ} \kn{ಮಾಲಿನ್ಯ} \kn{ಹೆಚ್ಚಾಗಿದೆ}.
\end{quote}

\begin{quote}
\emph{nagaradalli vāhanagaḷinda gāḷiya mālinya heccāgide.}
\end{quote}

\begin{quote}
In the city air pollution from vehicles has gone up.
\end{quote}

\kn{ವಾಯು} \kn{ಮಾಲಿನ್ಯ}, air pollution; \kn{ಜಲ} \kn{ಮಾಲಿನ್ಯ}, water pollution.

\begin{cousinweb}
\textbf{\kn{ಮಾಲಿನ್ಯ}} (\emph{mālinya}) is borrowed from Sanskrit \emph{mālinya}, \textquotedblleft{}foulness, dirtiness, impurity\textquotedblright{} (Monier-Williams), from \emph{malina}, \textquotedblleft{}dirty\textquotedblright{}.
\end{cousinweb}

\begin{grammarlens}[title={What you've built}]
One more word for getting around.
\end{grammarlens}

\subsection*{Your turn}
\begin{itemize}
\raggedright
  \item \emph{Say it:} \emph{mālinya}
  \item \emph{Say it:} \emph{mālinya}, once more
  \item \emph{Say it:} \emph{nagaradalli vāhanagaḷinda gāḷiya mālinya heccāgide.} — In the city air pollution from vehicles has gone up.
  \item \emph{Recall:} say the Kannada for a helpline, then the Kannada for the environment, then say \emph{mālinya} again
\end{itemize}

\begin{tcolorbox}[breakable,colback=teal!4,colframe=teal!35!black,title={Before you move on}]
What does \kn{ಮಾಲಿನ್ಯ} mean? (\textquotedblleft{}Pollution\textquotedblright{}.) Say it once more.
\end{tcolorbox}
\section[\texorpdfstring{\kn{ಕಲುಷಿತಗೊಳಿಸು} (kaluṣitagoḷisu)}{ಕಲುಷಿತಗೊಳಿಸು (kaluṣitagoḷisu)}]{\kn{ಕಲುಷಿತಗೊಳಿಸು} (kaluṣitagoḷisu) — to pollute; to contaminate}

\label{lesson:KA-C473-kalusitagolisu}



\begin{quote}
Before the new one: say the Kannada for the environment, then the Kannada for pollution.
\end{quote}

\subsection*{You'll want to know: \kn{ಕಲುಷಿತಗೊಳಿಸು}}
\textbf{\kn{ಕಲುಷಿತಗೊಳಿಸು}} — \emph{kaluṣitagoḷisu} — \textquotedblleft{}to pollute; to contaminate\textquotedblright{}.

\begin{quote}
\kn{ಕಾರ್ಖಾನೆ} \kn{ನದಿಯ} \kn{ನೀರನ್ನು} \kn{ಕಲುಷಿತಗೊಳಿಸಿತು}.
\end{quote}

\begin{quote}
\emph{kārkhāne nadiya nīrannu kaluṣitagoḷisitu.}
\end{quote}

\begin{quote}
The factory polluted the river water.
\end{quote}

Past (I …): \kn{ಕಲುಷಿತಗೊಳಿಸಿದೆ}, \emph{kaluṣitagoḷiside}.

\kn{ಕಲುಷಿತ} \kn{ನೀರು}, polluted water.

\begin{cousinweb}
\textbf{\kn{ಕಲುಷಿತಗೊಳಿಸು}} (\emph{kaluṣitagoḷisu}) is built from Sanskrit \emph{kaluṣita}, \textquotedblleft{}foul, impure\textquotedblright{} (Monier-Williams), and the Kannada verb-maker -\emph{goḷisu}: to make foul.
\end{cousinweb}

\begin{grammarlens}[title={What you've built}]
One more word for getting around.
\end{grammarlens}

\subsection*{Your turn}
\begin{itemize}
\raggedright
  \item \emph{Say it:} \emph{kaluṣitagoḷisu}
  \item \emph{Say it:} \emph{kaluṣitagoḷisu}, once more
  \item \emph{Say it:} \emph{kārkhāne nadiya nīrannu kaluṣitagoḷisitu.} — The factory polluted the river water.
  \item \emph{Recall:} say the Kannada for the environment, then the Kannada for pollution, then say \emph{kaluṣitagoḷisu} again
\end{itemize}

\begin{tcolorbox}[breakable,colback=teal!4,colframe=teal!35!black,title={Before you move on}]
What does \kn{ಕಲುಷಿತಗೊಳಿಸು} mean? (\textquotedblleft{}To pollute; to contaminate\textquotedblright{}.) Say it once more.
\end{tcolorbox}
\section[\texorpdfstring{\kn{ಮಸಿ} (masi)}{ಮಸಿ (masi)}]{\kn{ಮಸಿ} (masi) — soot}

\label{lesson:KA-C473-masi}



\begin{quote}
Before the new one: say the Kannada for pollution, then the Kannada for to pollute.
\end{quote}

\subsection*{You'll want to know: \kn{ಮಸಿ}}
\textbf{\kn{ಮಸಿ}} — \emph{masi} — \textquotedblleft{}soot\textquotedblright{}.

\begin{quote}
\kn{ಹಳೆಯ} \kn{ಬಸ್ಸಿನ} \kn{ಹೊಗೆಯಿಂದ} \kn{ಗೋಡೆಗೆ} \kn{ಮಸಿ} \kn{ಹಿಡಿಯಿತು}.
\end{quote}

\begin{quote}
\emph{haḷeya bassina hogeyinda gōḍege masi hiḍiyitu.}
\end{quote}

\begin{quote}
The smoke of an old bus left soot on the wall.
\end{quote}

In older Kannada it was also ink.

\begin{cousinweb}
\textbf{\kn{ಮಸಿ}} (\emph{masi}) is native Dravidian: DEDR 5101 lists Kannada \emph{masi}, \textquotedblleft{}soot, lampblack, ink\textquotedblright{}, beside Tamil \emph{mai}. Sanskrit \emph{maṣi}, \emph{masi}, \textquotedblleft{}lampblack, ink\textquotedblright{}, was itself taken from Dravidian, as Turner notes (CDIAL 9920).
\end{cousinweb}

\begin{grammarlens}[title={What you've built}]
One more word for getting around.
\end{grammarlens}

\subsection*{Your turn}
\begin{itemize}
\raggedright
  \item \emph{Say it:} \emph{masi}
  \item \emph{Say it:} \emph{masi}, once more
  \item \emph{Say it:} \emph{haḷeya bassina hogeyinda gōḍege masi hiḍiyitu.} — The smoke of an old bus left soot on the wall.
  \item \emph{Recall:} say the Kannada for pollution, then the Kannada for to pollute, then say \emph{masi} again
\end{itemize}

\begin{tcolorbox}[breakable,colback=teal!4,colframe=teal!35!black,title={Before you move on}]
What does \kn{ಮಸಿ} mean? (\textquotedblleft{}Soot\textquotedblright{}.) Say it once more.
\end{tcolorbox}
\section[\texorpdfstring{\kn{ಇಂಗಾಲ} (iṅgāla)}{ಇಂಗಾಲ (iṅgāla)}]{\kn{ಇಂಗಾಲ} (iṅgāla) — carbon}

\label{lesson:KA-C473-ingala}



\begin{quote}
Before the new one: say the Kannada for to pollute, then the Kannada for soot.
\end{quote}

\subsection*{You'll want to know: \kn{ಇಂಗಾಲ}}
\textbf{\kn{ಇಂಗಾಲ}} — \emph{iṅgāla} — \textquotedblleft{}carbon\textquotedblright{}.

\begin{quote}
\kn{ಮರಗಳು} \kn{ಗಾಳಿಯಿಂದ} \kn{ಇಂಗಾಲವನ್ನು} \kn{ಹೀರುತ್ತವೆ}.
\end{quote}

\begin{quote}
\emph{maragaḷu gāḷiyinda iṅgālavannu hīruttave.}
\end{quote}

\begin{quote}
Trees take in carbon from the air.
\end{quote}

In everyday Kannada a burnt-out coal; in science, carbon.

\begin{cousinweb}
\textbf{\kn{ಇಂಗಾಲ}} (\emph{iṅgāla}) comes from Sanskrit \emph{aṅgāra}, \textquotedblleft{}glowing charcoal\textquotedblright{}, through the variant \emph{iṅgāla} (CDIAL 125).
\end{cousinweb}

\begin{grammarlens}[title={What you've built}]
One more word for getting around.
\end{grammarlens}

\subsection*{Your turn}
\begin{itemize}
\raggedright
  \item \emph{Say it:} \emph{iṅgāla}
  \item \emph{Say it:} \emph{iṅgāla}, once more
  \item \emph{Say it:} \emph{maragaḷu gāḷiyinda iṅgālavannu hīruttave.} — Trees take in carbon from the air.
  \item \emph{Recall:} say the Kannada for to pollute, then the Kannada for soot, then say \emph{iṅgāla} again
\end{itemize}

\begin{tcolorbox}[breakable,colback=teal!4,colframe=teal!35!black,title={Before you move on}]
What does \kn{ಇಂಗಾಲ} mean? (\textquotedblleft{}Carbon\textquotedblright{}.) Say it once more.
\end{tcolorbox}
